% GENERATED FILE. Edit canonical lessons, then run npm run generate:books.
% canonical-source-hash: fnv1a64:32d77f91cb0b719d
% canonical-lessons: IT-C502-campeggiare, IT-C502-torcia, IT-C502-fornelletto, IT-C502-escursione, IT-C502-esplorare, IT-R502-ripasso-495-502

\chapter{Sotto le stelle}
\label{ch:it-sotto-le-stelle}
\hlchaptermodality{502}

\begin{chapteropening}
\textbf{By the end of this chapter:} \emph{I can talk about camping and exploring outdoors.}

Complete the last lesson of chapter 502: I can talk about camping and exploring outdoors.
\end{chapteropening}

\section[campeggiare]{campeggiare — to camp}

\label{lesson:IT-C502-campeggiare}



\begin{quote}
Before the new one: say the Italian for to bite, then the Italian for water bottle.
\end{quote}

\subsection*{You'll want to know: campeggiare}
\textbf{campeggiare} — \textquotedblleft{}to camp\textquotedblright{}.

\begin{quote}
D'estate campeggiamo vicino al fiume.
\end{quote}

\begin{quote}
\emph{In summer we camp near the river.}
\end{quote}

Built on \emph{il campeggio} \textquotedblleft{}campsite, camping\textquotedblright{}.

\begin{cousinweb}
\emph{campeggiare} is built on \emph{campo} \textquotedblleft{}field\textquotedblright{}, inherited from Latin \emph{campus} \textquotedblleft{}a plain, a field\textquotedblright{} (Lewis \& Short).
\end{cousinweb}

\begin{grammarlens}[title={What you've built}]
One more word for telling, explaining and describing at B1.
\end{grammarlens}

\subsection*{Your turn}
\begin{itemize}
\raggedright
  \item \emph{Say it:} \emph{campeggiare}
  \item \emph{Say it:} \emph{campeggiare}, once more
  \item \emph{Say it:} a sentence of your own with \emph{campeggiare}
  \item \emph{Recall:} say the Italian for to bite, then the Italian for water bottle, then say \emph{campeggiare} again
\end{itemize}

\begin{tcolorbox}[breakable,colback=teal!4,colframe=teal!35!black,title={Before you move on}]
What does campeggiare mean? (\textquotedblleft{}To camp\textquotedblright{}.) Say it once more.
\end{tcolorbox}
\section[la torcia]{la torcia — torch, flashlight}

\label{lesson:IT-C502-torcia}



\begin{quote}
Before the new one: say the Italian for water bottle, then the Italian for to camp.
\end{quote}

\subsection*{You'll want to know: la torcia}
\textbf{la torcia} — \textquotedblleft{}torch, flashlight\textquotedblright{}.

\begin{quote}
Di notte usiamo la torcia per trovare la tenda.
\end{quote}

\begin{quote}
\emph{At night we use the torch to find the tent.}
\end{quote}

In speech also \emph{la pila}.

\begin{cousinweb}
\emph{torcia} is usually traced to Vulgar Latin \emph{torca} \textquotedblleft{}something twisted\textquotedblright{}, from Latin \emph{torquēre} \textquotedblleft{}to twist\textquotedblright{} (Lewis \& Short): a torch was twisted straw or rope.
\end{cousinweb}

\begin{grammarlens}[title={What you've built}]
One more word for telling, explaining and describing at B1.
\end{grammarlens}

\subsection*{Your turn}
\begin{itemize}
\raggedright
  \item \emph{Say it:} \emph{la torcia}
  \item \emph{Say it:} \emph{la torcia}, once more
  \item \emph{Say it:} a sentence of your own with \emph{la torcia}
  \item \emph{Recall:} say the Italian for water bottle, then the Italian for to camp, then say \emph{la torcia} again
\end{itemize}

\begin{tcolorbox}[breakable,colback=teal!4,colframe=teal!35!black,title={Before you move on}]
What does la torcia mean? (\textquotedblleft{}Torch, flashlight\textquotedblright{}.) Say it once more.
\end{tcolorbox}
\section[il fornelletto]{il fornelletto — camping stove}

\label{lesson:IT-C502-fornelletto}



\begin{quote}
Before the new one: say the Italian for to camp, then the Italian for torch.
\end{quote}

\subsection*{You'll want to know: il fornelletto}
\textbf{il fornelletto} — \textquotedblleft{}camping stove\textquotedblright{}.

\begin{quote}
Abbiamo cucinato la pasta sul fornelletto.
\end{quote}

\begin{quote}
\emph{We cooked the pasta on the camping stove.}
\end{quote}

A little \emph{fornello} \textquotedblleft{}burner\textquotedblright{}.

\begin{cousinweb}
\emph{fornelletto} is a little form of \emph{fornello} \textquotedblleft{}burner\textquotedblright{}, built on \emph{forno} \textquotedblleft{}oven\textquotedblright{}, inherited from Latin \emph{furnus} \textquotedblleft{}an oven\textquotedblright{} (Lewis \& Short).
\end{cousinweb}

\begin{grammarlens}[title={What you've built}]
One more word for telling, explaining and describing at B1.
\end{grammarlens}

\subsection*{Your turn}
\begin{itemize}
\raggedright
  \item \emph{Say it:} \emph{il fornelletto}
  \item \emph{Say it:} \emph{il fornelletto}, once more
  \item \emph{Say it:} a sentence of your own with \emph{il fornelletto}
  \item \emph{Recall:} say the Italian for to camp, then the Italian for torch, then say \emph{il fornelletto} again
\end{itemize}

\begin{tcolorbox}[breakable,colback=teal!4,colframe=teal!35!black,title={Before you move on}]
What does il fornelletto mean? (\textquotedblleft{}Camping stove\textquotedblright{}.) Say it once more.
\end{tcolorbox}
\section[l'escursione]{l'escursione — trip, hike, excursion}

\label{lesson:IT-C502-escursione}



\begin{quote}
Before the new one: say the Italian for torch, then the Italian for camping stove.
\end{quote}

\subsection*{You'll want to know: l'escursione}
\textbf{l'escursione} — \textquotedblleft{}trip, hike, excursion\textquotedblright{}.

\begin{quote}
Domani facciamo un'escursione in montagna.
\end{quote}

\begin{quote}
\emph{Tomorrow we are going on a hike in the mountains.}
\end{quote}

\emph{Fare un'escursione} — to go on a trip or a hike.

\begin{cousinweb}
\emph{escursione} is a learned borrowing from Latin \emph{excursiō} \textquotedblleft{}a running out, a setting out\textquotedblright{} (EtymDB), from \emph{excurrere} \textquotedblleft{}to run out\textquotedblright{}.
\end{cousinweb}

\begin{grammarlens}[title={What you've built}]
One more word for telling, explaining and describing at B1.
\end{grammarlens}

\subsection*{Your turn}
\begin{itemize}
\raggedright
  \item \emph{Say it:} \emph{l'escursione}
  \item \emph{Say it:} \emph{l'escursione}, once more
  \item \emph{Say it:} a sentence of your own with \emph{l'escursione}
  \item \emph{Recall:} say the Italian for torch, then the Italian for camping stove, then say \emph{l'escursione} again
\end{itemize}

\begin{tcolorbox}[breakable,colback=teal!4,colframe=teal!35!black,title={Before you move on}]
What does l'escursione mean? (\textquotedblleft{}Trip, hike, excursion\textquotedblright{}.) Say it once more.
\end{tcolorbox}
\section[esplorare]{esplorare — to explore}

\label{lesson:IT-C502-esplorare}



\begin{quote}
Before the new one: say the Italian for camping stove, then the Italian for trip.
\end{quote}

\subsection*{You'll want to know: esplorare}
\textbf{esplorare} — \textquotedblleft{}to explore\textquotedblright{}.

\begin{quote}
I bambini vogliono esplorare il bosco.
\end{quote}

\begin{quote}
\emph{The children want to explore the woods.}
\end{quote}

The person is \emph{l'esploratore}.

\begin{cousinweb}
\emph{esplorare} is a learned borrowing from Latin \emph{explōrāre} \textquotedblleft{}to search out, to examine\textquotedblright{} (EtymDB; Lewis \& Short).
\end{cousinweb}

\begin{grammarlens}[title={What you've built}]
One more word for telling, explaining and describing at B1.
\end{grammarlens}

\subsection*{Your turn}
\begin{itemize}
\raggedright
  \item \emph{Say it:} \emph{esplorare}
  \item \emph{Say it:} \emph{esplorare}, once more
  \item \emph{Say it:} a sentence of your own with \emph{esplorare}
  \item \emph{Recall:} say the Italian for camping stove, then the Italian for trip, then say \emph{esplorare} again
\end{itemize}

\begin{tcolorbox}[breakable,colback=teal!4,colframe=teal!35!black,title={Before you move on}]
What does esplorare mean? (\textquotedblleft{}To explore\textquotedblright{}.) Say it once more.
\end{tcolorbox}
\section[(dialogue)]{Ripasso dei capitoli 495–502}

\label{lesson:IT-R502-ripasso-495-502}



\begin{quote}
Say the Italian for to adopt and adoption.
\end{quote}

\subsection*{Your turn}
Say these aloud:

\begin{itemize}
\raggedright
  \item to adopt, adoption, stepfather, stepmother, half-brother
  \item nappy, pushchair, to breastfeed, to rock, to crawl
  \item to bring up, to spoil, spoilt, to scold, punishment
  \item nanny, paediatrician, to care for, to look after, to obey
  \item argument, to make up, to betray, faithfulness, to confide
  \item to sew, needle, to embroider, crochet, to darn
  \item to fish, fishing rod, bait, to bite, water bottle
  \item to camp, torch, camping stove, trip, to explore
\end{itemize}

\begin{tcolorbox}[breakable,colback=teal!4,colframe=teal!35!black,title={Before you move on}]
To adopt? (\textbf{adottare}.) Adoption? (\textbf{l'adozione}.) Stepfather? (\textbf{il patrigno}.) Stepmother? (\textbf{la matrigna}.) Half-brother? (\textbf{il fratellastro}.) Nappy? (\textbf{il pannolino}.) Pushchair? (\textbf{il passeggino}.) To breastfeed? (\textbf{allattare}.) To rock? (\textbf{cullare}.) To crawl? (\textbf{gattonare}.) To bring up? (\textbf{educare}.) To spoil? (\textbf{viziare}.) Spoilt? (\textbf{viziato}.) To scold? (\textbf{rimproverare}.) Punishment? (\textbf{il castigo}.) Nanny? (\textbf{la tata}.) Paediatrician? (\textbf{il pediatra}.) To care for? (\textbf{accudire}.) To look after? (\textbf{badare}.) To obey? (\textbf{ubbidire}.) Argument? (\textbf{il litigio}.) To make up? (\textbf{riconciliarsi}.) To betray? (\textbf{tradire}.) Faithfulness? (\textbf{la fedeltà}.) To confide? (\textbf{confidarsi}.) To sew? (\textbf{cucire}.) Needle? (\textbf{l'ago}.) To embroider? (\textbf{ricamare}.) Crochet? (\textbf{l'uncinetto}.) To darn? (\textbf{rammendare}.) To fish? (\textbf{pescare}.) Fishing rod? (\textbf{la canna da pesca}.) Bait? (\textbf{l'esca}.) To bite? (\textbf{abboccare}.) Water bottle? (\textbf{la borraccia}.) To camp? (\textbf{campeggiare}.) Torch? (\textbf{la torcia}.) Camping stove? (\textbf{il fornelletto}.) Trip? (\textbf{l'escursione}.) To explore? (\textbf{esplorare}.)
\end{tcolorbox}
